%% notes-tex/010-additive-baselines/sections/1-question.tex

\section{The question\secstatus{todo}}

Tran et al.\ introduced MULTI-evolve in \emph{Science} in 2026
(\texttt{10.1126/science.aea1820}), a protein engineering framework whose fully
connected networks were said to ``learn the epistatic landscape'' from single
and double mutants and to identify synergistic combinations at higher order.
Visani, Verma and DeWitt replied in a preprint,
\emph{Additive baselines furnish no evidence for epistasis learning by
MULTI-evolve} (bioRxiv \texttt{2026.04.23.719915}), which does one thing: it
fits the null that the original paper never fit.

An \term{additive model} over mutations predicts a phenotype as a sum of
per-mutation terms. By construction it assigns exactly zero to every
interaction. It is therefore the natural null for any claim that a model has
learned interaction, in the same way that an intercept-only model is the null
for a claim that a model has learned anything.

Their result is that the network's predictions across the combinatorial variant
space correlate with a ridge-regularized additive model's at $r > 0.999$ for all
three engineering campaigns, that the variants selected for construction are
ranked nearly identically by both, and that on held-out test data the network
does not significantly outperform either the additive model or a zero-hidden-layer
version of its own architecture. The engineering result is real; the attribution
to epistasis learning is not supported.

The analogous question for experiment 010 is whether the equivariant cell graph
transformer's held-out performance on the trigenic interaction score survives
the same test. Traditional machine-learning baselines exist in this repository
for the mixed digenic and trigenic set under \texttt{experiments/002-dmi-tmi}
and for fitness under \texttt{experiments/smf-dmf-tmf-001}. None had been fit on
the 010 build, which is triples only.

\notesec{One difference in the setup, and it cuts against the null}
In the protein setting the quantity predicted is the multimutant phenotype
itself, which per-mutation effects explain directly. In 010 the quantity
predicted is already a residual. The trigenic interaction score subtracts the
expectation built from single and double mutant fitness before the model ever
sees it, so the additive part of fitness has been removed by construction. An
additive-in-single-gene-effects model is therefore a harder null here than it is
there, and its performance is more surprising, not less.

That argument has one loophole, and quantifying it turns out to matter more than
anything else in this document. It is taken up in
Section~\ref{sec:results} as baseline B4.
